\documentclass{jsarticle}
\usepackage[dvipdfmx, final]{graphicx}
\usepackage{amsmath}
\usepackage{amssymb}
\usepackage{chemfig}
\begin{document}
\title{重力波分解での電磁場発電所}
\author{Masaaki Yamaguchi}
\date{}
\maketitle


        熱力学の第１法則は、熱エネルギー保存則、
   熱力学の第２法則は、エントロピー増大則、
   これより、第１法則より、摩擦熱があると、永久機関はつくられない。
   マクスウェルの方程式より、QEDから、重力を分解すると、電磁場が発生する。
   この電磁場から発生する電磁気力を分解すると、弱い力と強い力ができる。

   第１法則から、逆の統合はない。
   すでに、統一場理論は、ファインマン博士とアインシュタイン博士により、
   見つかっている。
   ガンマ関数の大域的部分積分多様体が、この統一場理論になっている。
   $$ T^{'} = \int \Gamma(\gamma)^{'}dx_m $$
   重力であるダランベルシアンの中のガンマ関数を共変微分すると、
   これは、実は大域的微分であり、共変微分とは、全然違う計算方法である。
   この重力の積分の中で微分するとは、上の説明を鑑みるとわかる。
   結果は、同じなれど、統一場理論になっている。　
   電磁場を弱い力と強い力に統合できるが、それが電弱相互作用であるが、
   熱力学の第１法則より、重力とは、この法則の逆より、本当は、重力を
   分解していくと、宇宙の真理の統一場理論が、重力だけで事足りるのがわかる。
   逆の統合はないとおもう。カタストロフィー理論から、形態形成場理論も
   あり、自己組織化現象もあるので、自信はない。
   時間の流れが未来から過去より、この時間が、ガンマ関数の大域的部分積分
   多様体より、あるかもしれない。
   逆を探すより、重力を分解していく統一場理論の方がすんなり言える。
   情報科学の世界では、逆がある。
   コップが割れると、破片になり、逆は、コップがくっつかないともとに戻せない。
   素粒子同士をくっつかせているのが、中間子力であり、
   世界なにがあるかわからない。



       原子力発電所の代わりは、重力波をプリズム体でQED分解すると、
   電磁場が発生して、それを更に分解すると、弱い力と強い力に最終的に
   分かれる。この重力波を電磁場にするのを、電磁場発電所として開発すると、
   何兆億という、一代企業も夢ではないとおもう。竹内薫先生も、
   ブラックホールでのゴミでの、燃やしたときの熱エネルギーで、
   発電所ができることを、ファンに言っている。
   これは、宇宙空間でするので、地球温暖化は関係ない。
   弱い力と強い力は、核開発に使われると、湯川秀樹先生が中間子の
   媒介エネルギーで、断固反対と言っている。
   やはり、重力を電磁場にするのが一番いいとおもう。




        このプリズム体は、仕組みは、サーストン・ペレルマン多様体をヒルベルト
   空間の中に、双対被覆でのグラスマン多様体として、組み合わせ多様体での
   構成とするのを、部品として使うようにして、これをLSIで熱感知器として、
   Jones多項式で、この回路のFPGAに組み込むとできるとおもう。
   最適化は、どの空間でも、熱力学第２法則で必要である。
   人体には、反重力での磁場発生で解決している。
   電磁場を回転体として、使うと、反重力が発生するのをNECの南先生が、
   航空機で研究している。
   これは、東北大学の早坂先生も考えている。　


\end{document}
